% Laboratorio 7 de Automatizacion - Informe IEEE editable.
% Informe en progreso: ensayo ABIERTO real incorporado; observadores físicos aún pendientes.
% Guia oficial con encabezado erroneo "Laboratorio 6".
\documentclass[conference,10pt]{IEEEtran}
\usepackage[T1]{fontenc}
\usepackage[utf8]{inputenc}
\usepackage[spanish,es-nodecimaldot]{babel}
\usepackage{amsmath,amssymb,mathtools}
\usepackage{booktabs,tabularx,array}
\usepackage{graphicx}
\usepackage{cite}
\usepackage{siunitx}
\usepackage{url}
\usepackage[hidelinks]{hyperref}
\sisetup{output-decimal-marker={,},per-mode=symbol}

\title{Diseño e Implementación de Controladores Digitales\\
con Estimadores de Predicción y Estado Actual}
\author{\IEEEauthorblockN{[Nombre del estudiante]}\\
\IEEEauthorblockA{Carrera de Ingeniería Electrónica\\
Universidad Católica ``Nuestra Señora de la Asunción''\\
Asignatura: Automatización --- Laboratorio N.\textsuperscript{o} 7\\
Docente: PhD. Enrique A. Vargas Cabral}}

\begin{document}
\maketitle
\begin{abstract}
Se presenta la metodología de diseño e implementación de controladores digitales
por realimentación de estados estimados para una planta analógica de segundo orden.
Se consideran un estimador de predicción y un estimador de estado actual, con dos
ubicaciones alternativas de polos del observador. El desarrollo se organiza desde
el modelo físico y la selección del tiempo de muestreo hasta la discretización,
la asignación de polos, la simulación en MATLAB y la implementación en PSoC 5LP.
La cadena de adquisición y actuación se documenta a partir de los antecedentes
del Laboratorio 5 y se establece un protocolo de revalidación con el osciloscopio
Rigol DS1202Z-E. Se obtuvo el modelo nominal del esquema oficial, se discretizó
a \SI{2}{\milli\second} y se verificó controlabilidad y observabilidad.
Se validó físicamente la planta en lazo abierto: se registraron
simultáneamente el pulso en J5 y la salida en J7 mediante el
Rigol DS1202Z-E y 256 muestras por UART en el PSoC.
La implementación y evaluación física de los dos observadores sigue pendiente.
\end{abstract}
\begin{IEEEkeywords}
control digital, estimación de estados, observador de predicción,
observador actual, ubicación de polos, PSoC 5LP, MATLAB.
\end{IEEEkeywords}

\section{Introducción}
La realimentación completa de estados permite modificar los polos del sistema,
pero requiere conocer cada variable del vector de estado. Cuando sólo es posible
medir la salida de la planta, el diseño digital puede incluir un observador
capaz de reconstruir estados no medidos a partir del modelo, el mando y la
señal disponible \cite{samanta2024,fadali2020,astrom2020}.

La guía de esta práctica exige diseñar e implementar dos estructuras: un
estimador de predicción y un estimador de estado actual. Ambas se aplicarán a
una planta analógica de segundo orden. A diferencia del laboratorio previo,
la realimentación debe efectuarse \emph{únicamente sobre los estados estimados};
la medición física usada por el algoritmo será la salida del sistema.

La experiencia acumulada en el Laboratorio 5 proporciona una infraestructura
experimental de PSoC 5LP, MATLAB y adquisición con el Rigol DS1202Z-E. Sin
embargo, la calibración de esa infraestructura y el modelo de la nueva práctica
requieren verificación antes de cerrar el lazo. Esta distinción permite preservar
trazabilidad sin trasladar valores antiguos como resultados actuales.

\section{Objetivos}
\subsection{Objetivo general}
Diseñar e implementar controladores digitales para una planta analógica mediante
realimentación de estados estimados, utilizando observadores de predicción y de
estado actual, de acuerdo con la guía oficial del Laboratorio 7.

\subsection{Objetivos específicos}
\begin{enumerate}
\item Seleccionar y justificar el tiempo de muestreo antes de discretizar.
\item Obtener el modelo en espacio de estados y comprobar controlabilidad y observabilidad desde la salida medida.
\item Diseñar una ley de control común por ubicación de polos.
\item Obtener los observadores de predicción y actual; evaluar dos ubicaciones de sus polos.
\item Comparar las respuestas, errores de estimación y esfuerzos de control.
\item Implementar ambos esquemas en PSoC y confrontar los ensayos con MATLAB.
\end{enumerate}

\section{Plataforma experimental y antecedentes}
\subsection{Instrumentación}
La plataforma comprende un PSoC 5LP con ADC, DAC y temporizador de muestreo;
planta analógica con acondicionamiento de entrada y salida; fuente de
alimentación; PC con MATLAB; y osciloscopio \emph{Rigol DS1202Z-E}.
La instrumentación previa empleó USB-TMC/SCPI, sondas \(10\times\),
acoplamiento DC, telemetría UART a \SI{57600}{bit\per\second} y adquisición
sincronizada. Estos antecedentes constituyen una referencia operativa, no una
calibración vigente del nuevo montaje.

\subsection{Calibración histórica}
La Tabla~\ref{tab:calib_historica} reúne parámetros documentados en los
archivos de configuración y en el ajuste experimental del Laboratorio 5.
La adquisición del Lab7 comenzará con una \textbf{revalidación} de cero, rango,
pendientes y ruido a partir de su propio cableado.

\begin{table}[t]
\centering
\caption{Parámetros documentados del Laboratorio 5 (no revalidados en Lab7).}
\label{tab:calib_historica}
\small
\begin{tabular}{@{}ll@{}}
\toprule
Parámetro histórico & Valor documentado \\
\midrule
Offset del ADC & \SI{2.353802067}{\volt} \\
Escala ADC/planta & \SI{0.210251951}{\volt\per\volt} \\
Ganancia de DAC (+) & 3,842612 \\
Ganancia de DAC (-) & 2,640105691 \\
Offset físico J5 & \SI{17.209}{\milli\volt} \\
Pendiente J5/mando aplicado & 0,981710 \\
RMSE del ajuste de J7 & \SI{3.374}{\milli\volt} \\
\bottomrule
\end{tabular}
\end{table}

En las comparaciones se conservarán las tensiones crudas de los dos canales,
la telemetría digital, los valores calibrados, la referencia física, la
señal de control realmente aplicada y los metadatos de disparo del osciloscopio.
Esto evita confundir coincidencias visuales de amplitud con validaciones
cuantitativas de la respuesta del sistema.

\section{Modelado y discretización}
\subsection{Modelo continuo}
A partir del diagrama analógico de la guía oficial, con definición explícita
de los nodos que constituyen el vector de estados, se construirá
\begin{align}
 \dot{\mathbf{x}}(t) &= F\mathbf{x}(t)+G u(t),\\
 y(t)&=H\mathbf{x}(t)+J u(t).
\end{align}
Para una planta de segundo orden y una sola salida medida,
\(F\in\mathbb R^{2\times2}\), \(G\in\mathbb R^{2\times1}\),
\(H\in\mathbb R^{1\times2}\) y \(J\in\mathbb R\).
Las matrices numéricas, los valores R--C y las unidades se incorporarán
cuando se confirme el circuito físico; no se copiarán de un informe ajeno.

\subsection{Selección justificada del muestreo}
La selección de \(T_s\) precede al cálculo de \(A\) y \(B\). Se estudiarán la
frecuencia característica más rápida, el ancho de banda de diseño, la
relación de muestreo exigida en clase y las cuentas enteras del temporizador
\cite{samanta2024,fadali2020}. La frecuencia efectiva será
\begin{equation}
 f_s = \frac{1}{T_s},\qquad \omega_s=\frac{2\pi}{T_s}.
\end{equation}
El valor definitivo de \(T_s\) queda \emph{pendiente de verificación};
no se adopta por defecto el de los laboratorios anteriores.

\subsection{Modelo discreto con retenedor de orden cero}
Para mando constante entre muestras (ZOH),
\begin{align}
 A &= e^{FT_s}, & B &= \int_{0}^{T_s} e^{F\tau}G\,d\tau,\\
 C &= H, & D &= J. 
\end{align}
De ese modo,
\begin{equation}\label{eq:planta_discreta}
\mathbf{x}[k+1]=A\mathbf{x}[k]+Bu[k],\quad y[k]=C\mathbf{x}[k]+Du[k].
\end{equation}
El cálculo se verificará con \texttt{c2d(ss(F,G,H,J),Ts,'zoh')}.

\subsection{Controlabilidad y observabilidad}
Antes de asignar polos se comprobarán
\begin{equation}
 \mathcal C=[B\;AB],\qquad
 \mathcal O=\begin{bmatrix} C\\CA\end{bmatrix},
\end{equation}
y se exigirá \(\operatorname{rango}(\mathcal C)=
\operatorname{rango}(\mathcal O)=2\). La observabilidad debe verificarse
\emph{desde el único canal de salida} disponible en el experimento.

\section{Diseño del controlador y estimadores}
\subsection{Realimentación de estados estimados}
La ley de control para ambos esquemas será
\begin{equation}\label{eq:control}
 u[k]=-K\widehat{\mathbf{x}}[k]+N_r r[k].
\end{equation}
Los polos del controlador se asignarán mediante \(K=\operatorname{place}
(A,B,p_c)\), verificando los valores propios de \(A-BK\). En el caso
SISO, \(D=0\) y seguimiento de referencia constante, se utilizará
\begin{equation}
 N_r=\left[C(I-A+BK)^{-1}B\right]^{-1}.
\end{equation}
Los cálculos no suponen que el comportamiento con saturación sea exactamente
el del modelo lineal sin restricciones.

\subsection{Estimador de predicción}
La innovación se calcula con la salida disponible \(y[k]\), pero la
estimación corregida se utiliza en el siguiente instante:
\begin{equation}\label{eq:predictor}
\widehat{\mathbf{x}}[k+1]=A\widehat{\mathbf{x}}[k]+Bu[k]
 +L_p\big(y[k]-C\widehat{\mathbf{x}}[k]\big).
\end{equation}
Con \(D=0\), la dinámica de error es
\begin{equation}
 \mathbf e[k+1]=(A-L_pC)\mathbf e[k].
\end{equation}
Por dualidad, \(L_p=\operatorname{place}(A^T,C^T,p_e)^T\).

\subsection{Estimador de estado actual}
El observador actual distingue predicción y corrección en la misma muestra:
\begin{align}
\mathbf z[k]&=A\widehat{\mathbf{x}}[k-1]+Bu[k-1],\\
\widehat{\mathbf{x}}[k]&=\mathbf z[k]+L_a\big(y[k]-C\mathbf z[k]\big).
\end{align}
La ley de control~\eqref{eq:control} utiliza la estimación corregida.
La matriz de transición del error corregido es
\begin{equation}
 (I-L_a C)A=A-L_a C A.
\end{equation}
En consecuencia, para asignar sus polos se empleará
\(L_a=\operatorname{place}(A^T,(CA)^T,p_e)^T\), comprobando previamente
observabilidad y rango adecuados. Se evitará equiparar directamente
\(L_p\) con \(L_a\), pues representan arquitecturas diferentes.

\subsection{Experimentos de ubicación de polos}
Se mantendrá el mismo \(K\) en las cuatro comparaciones:
(predicción, E1), (predicción, E2), (actual, E1) y (actual, E2).
Las dos parejas del observador deberán situarse dentro del círculo unidad
y responder a un criterio de velocidad y sensibilidad al ruido.
La elección definitiva de polos y los parámetros numéricos se registrarán
en las tablas de diseño, no se fijan en esta versión inicial.

\section{Metodología de simulación e implementación}
\subsection{Simulación en MATLAB}
El programa \texttt{lab7\_principal.m} compara cuatro combinaciones de
observador utilizando la misma planta, referencia y condiciones iniciales.
Se exportarán series temporales y las matrices diseñadas. Se analizarán:
la respuesta \(y[k]\), el esfuerzo \(u[k]\), la innovación, las
trayectorias estimadas, el sobrepasamiento, el tiempo de establecimiento,
el error estacionario y la activación de límites de actuación. El error de
estimación completo \(\mathbf x-\widehat{\mathbf x}\) es observable en
simulación; en el banco debe distinguirse de las mediciones físicas
disponibles.

\subsection{Arquitectura propuesta para el PSoC}
En cada evento de muestreo se adquirirá la salida de la planta mediante el
ADC, se convertirá a unidades físicas y se actualizará el estimador
seleccionado. Después se calculará la ley~\eqref{eq:control}, se limitará
el mando a la región validada y se enviará al DAC. La UART almacenará
referencia, salida, estados estimados, innovación, mando solicitado,
mando aplicado, banderas y tiempos de ejecución. La temporización se
verificará con un pin de diagnóstico y el Rigol DS1202Z-E.

El programa C inicial contiene \texttt{lab7\_estimadores.c} y
\texttt{lab7\_estimadores.h}. Es el núcleo matemático y
\emph{todavía no constituye firmware de PSoC Creator ejecutable};
resta integrar el diseño de periféricos, ISR, conversión ADC/DAC y
protecciones reales.

\subsection{Protocolo de medición}
Antes del lazo cerrado se comprobarán la continuidad de los puntos de
prueba, las masas, el ajuste de sondas \(10\times\), la salida nula,
el rango bipolar del DAC, la calibración del ADC, el periodo muestreado
y el tiempo máximo de ejecución. Se guardarán simultáneamente ondas
crudas del Rigol, registros PSoC y metadatos, con una única amplitud de
referencia física por ensayo y sin correcciones tácitas.

\section{Resultados y discusión}
\subsection{Alcance de las simulaciones}
Se ejecutaron en MATLAB R2026b las cuatro configuraciones nominales
(predicción y estado actual; ubicaciones E1 y E2) con
$T_s=\SI{2}{\milli\second}$ y referencia de \SI{50}{\milli\volt}
aplicada en $t=\SI{20}{\milli\second}$.
La salida $y[k]$ es el único estado medido que se realimentaría físicamente.
Los valores de esta sección son \emph{exclusivamente simulaciones}, no
mediciones en PSoC. Las gráficas E2 permanecen pendientes de discusión.

\begin{table}[t]
\centering
\caption{Estado de los casos del Laboratorio 7.}
\label{tab:ensayos}
\small
\begin{tabular}{@{}llll@{}}
\toprule
Observador & Polos & MATLAB & PSoC \\
\midrule
Predicción & E1 & Analizado & Pendiente \\
Predicción & E2 & Ejecutado & Pendiente \\
Actual & E1 & Analizado & Pendiente \\
Actual & E2 & Ejecutado & Pendiente \\
\bottomrule
\end{tabular}
\end{table}

\subsection{Comparación nominal entre observadores E1}
Para ambos controladores se utilizaron los polos de lazo cerrado
$0.8\pm j0.25$, los polos de observación E1 $0.4\pm j0.25$
y las mismas condiciones iniciales: $x[0]=[0\quad0]^\top$,
$\hat{x}[0]=[0.02\quad-0.01]^\top$. Las métricas de la
Tabla~\ref{tab:res_e1} se calcularon reproduciendo numéricamente el
algoritmo de las simulaciones MATLAB, y se contrastaron con las
gráficas aportadas del ensayo computacional.

\begin{table}[t]
\centering
\caption{Comparación E1: simulación nominal, no ensayo físico.}
\label{tab:res_e1}
\small
\begin{tabular}{@{}lrr@{}}
\toprule
Magnitud & Predictor & Actual \\
\midrule
Máximo de salida (\si{\milli\volt}) & 57.646 & 57.631 \\
Sobrepasamiento (\%) & 15.293 & 15.263 \\
Pico desde el escalón (\si{\milli\second}) & 22 & 22 \\
Establecimiento $\pm2\%$ (\si{\milli\second}) & 46 & 48 \\
$u[0]$ (\si{\milli\volt}) & 25.57 & 55.90 \\
$\max u$ (\si{\milli\volt}) & 97.98 & 97.88 \\
$\min u$ (\si{\milli\volt}) & -8.72 & -16.25 \\
$u_\infty$ (\si{\milli\volt}) & 49.09 & 49.09 \\
\bottomrule
\end{tabular}
\end{table}

El predictor presentó un máximo de \SI{8.20}{\milli\volt}
antes de aplicar la referencia, frente a \SI{15.22}{\milli\volt}
del estimador actual. Esta diferencia surge del error de inicialización:
en el estimador actual, la medición de la muestra $k$ modifica la
estimación antes de calcular $u[k]$; el predictor utiliza la innovación
para construir el estado estimado de la siguiente muestra.
Aunque el estimador actual utiliza información más reciente, no
proporciona por sí mismo un menor tiempo de establecimiento: en este
caso, ambos comportamientos son muy semejantes.

La norma de error inicial reportada por el estimador actual está
calculada \emph{después} de corregir su predicción con $y[0]$,
mientras que la del predictor es previa a esa corrección. Por ello
no debe compararse directamente la altura inicial de las curvas
como si representaran el mismo instante de cálculo.
En condiciones ideales, ambas arquitecturas presentan convergencia
del error; resta estudiar ruido, cuantización y retardos reales.

\subsection{Ensayos experimentales pendientes}
Para cada diseño deberán medirse los tiempos de establecimiento,
los sobrepasamientos, el error estacionario, el máximo y el RMS del
esfuerzo, y el efecto del ruido. La comparación MATLAB--banco requiere
datos exportados del Rigol DS1202Z-E y calibraciones ADC/DAC actuales.
No se atribuyen las métricas nominales al hardware físico.

\section{Conclusiones y trabajo pendiente}
La guía oficial establece un contraste experimental entre control con
estimador de predicción y con estimador actual. En esta etapa quedó
definida una metodología reproducible para modelar, verificar
controlabilidad/observabilidad, asignar polos, simular ambas estructuras
y documentar su implementación. Se conservó la trazabilidad de las
calibraciones históricas del Laboratorio 5 sin atribuirles vigencia
automática para el nuevo circuito.

\textbf{Conclusión final pendiente:} deberá redactarse a partir de
mediciones del Laboratorio 7 y del contraste cuantitativo con MATLAB,
indicando diferencias, causas técnicas, limitaciones y cumplimiento de
los objetivos. Esta redacción preliminar no sustituye la conclusión
experimental exigida por la cátedra.


\begin{thebibliography}{9}
\bibitem{samanta2024} B. Samanta, \emph{Introduction to Digital Control: An Integrated Approach}. Springer, 2024.
\bibitem{fadali2020} M. S. Fadali y A. Visioli, \emph{Digital Control Engineering: Analysis and Design}, 3.\textsuperscript{a} ed. Academic Press, 2020.
\bibitem{astrom2020} K. J. Åstr\"om y R. M. Murray, \emph{Feedback Systems: An Introduction for Scientists and Engineers}, 2.\textsuperscript{a} ed., versión electrónica v3.1.5, 2020.
\bibitem{franklin1998} G. F. Franklin, J. D. Powell y M. L. Workman, \emph{Digital Control of Dynamic Systems}, 3.\textsuperscript{a} ed. Addison-Wesley, 1998.
\bibitem{guiaLab7} Cátedra de Automatización, Universidad Católica Nuestra Señora de la Asunción, ``Utilización y evaluación de estimadores de estado de predicción y actual'', guía oficial de laboratorio (rótulo numérico de origen erróneo).
\end{thebibliography}

\end{document}
